% !TEX root = 00_main.tex
\chapter{序論}\label{chap:introduction}

\section{背景}\label{sec:background}
% 社会的背景 → 現状の問題 → 解決すべき課題 の順に書く．
% 主張には必ず根拠（参考文献）を付ける．例：\cite{sample-book}
ここに研究の背景を書く．
\todo{ここに統計データを追加する}
文献を引用する場合は\verb|\cite{キー}|を用いる\cite{sample-book}\cite{sample-web}．
書きかけの箇所には\verb|\todo{...}|でメモを残せる（提出前に\verb|01_macro.tex|の設定で一括非表示にできる）．

\section{目的}\label{sec:purpose}
% 本研究で何を明らかにする／何を作るのかを明確に書く．
ここに研究の目的を書く．

\section{本論文の構成}
本論文の構成は以下の通りである．
\chapref{related-work}では関連研究について述べ，本研究の位置づけを示す．
\chapref{method}では提案手法について述べる．
\chapref{implementation}では提案手法の実装について述べる．
\chapref{evaluation}では評価実験とその結果について述べる．
最後に\chapref{conclusion}で結論を述べる．
